% Page 053

\seite{53}

das Schuljahr 1915/16\ob
\margmark{Entlassen wurde kei\\chen. Die Schülerza\\Kinder sind katholi}%
n Schüler, aufgenommen wurden 5 Mäd\ohb
hl beträgt 32 – 13 Knaben, 19 Mädchen. Alle\ob
sch.

\pend
\abschnitt{Ernte}
\pstart

\margmark{Die Gemeinde erntet\\Gerste und Hafer ge\\Die Wiesenpreise wa\\Butter, die Preise\\Jahre sehr viel höh\\Preise 1,00 – 1,20}%
e in diesem Jahre eine mittlere Ernte. Weizen\ob
diehen gut. Die Kartoffelernte ist zufriedenstellend.\ob
ren den Sommer über sehr hoch, besonders die\ob
der Lebensmittel waren gegen den vorstehenden\ob
er. Butter und Eier kosteten 1,50 – 1,60 M, während Eier die\ob
M betrugen.

\pend
\abschnitt{Fortsetzung}
\pstart

\margmark{Während unsere Sold\\zur See mit bewunde\\des Feindes ertrage\\forthalten, wird au\\schwere Opfer ihres\\die Obligation geht\\längst eingelagen.\\hat die Arbeit viel\\helfen sich fleißig\\lich zu erweisen, u\\der Christlichen Bü\\die übrigen Klassen\\tagsschule im Somme\\in den Volksschulen\\Jahresbericht die S\\für Kriegswaisen. E\\ein Soldat ist seit}%
aten im Osten und Westen, zu Land und\ob
rungswürdiger Tapferkeit die Truppen des\ob
n, und so das Land von Feind und Hunger\ob
ch die Bevölkerung aufs neue, weitere\ob
 Vaterlandes zu bringen. Zunächst steht auf,\ob
 zurück. Die besten Rohstoffe sind\ob
Landes- und Kreisausschüsse, die eingesetzt haben,\ob
fach geleistet, fort. Kinder und alte Leute\ob
e Arbeit dem lieben Vaterlande nütz\ohb
nd in diesem Sinne werden die Kinder\ob
rgerschule fort vom Schuljahre befreit, außerdem\ob
, daß Anordnung der Behörden ein Halb\ohb
r halten. Zu wiederholten Male werden\ob
 frei von Krieg für den Vaterländischen\ob
chulkinder erfahren jede Vergünstigungen\ob
s sind fünf im Dorfe 4 Krieger gestellt,\ob
 einem Jahre vermißt. Die ganze Pfarrei
\pend
